% -----------------------------
% 執筆にあたっての留意点
% -----------------------------
%
% 「提案事項の詳細」について触れる章の名前は内容に依ります．単に「提案手法」「提案システム」とする場合もあれば，「XXを行う手法」と具体的な名前をつけることもあります．また，章を複数に分割する場合もあります．
%
% 提案手法のアイデアや挙動を説明するときは，必ず例を使いながら説明してください．そうすることで，読者が内容を理解しやすくなります．ここで使う例は「イントロダクション」で使った例と同じものを使うのが望ましいです．
%
%
% -----------------------------
% 補足（参考：how-to-xx/how-to-write-a-paper/1-学術論文の書き方.md 3.4節）
% -----------------------------
%
% ■ 例は使い捨てにせず，一つの例を論文全体で育ててください
% 上に書いた「イントロダクションと同じ例を使う」ことの狙いです．場面ごとに新しい例を出すと，
% 読者は新しい概念が出てくるたびに新しい状況を理解し直さなければなりません．
% 一つの例を，定義 → 手法 → 実行例 → 実験と一貫して使えば，読者は認知負荷を抑えたまま読み進められます．
%
% ■ 論文全体のおよそ1/4を読み終えるまでに，新規性のある技術的貢献を提示してください
% 8ページの論文なら2ページ目の終わりごろ，12ページの論文なら3ページ目あたりが目安です．
% そこまで読んでも「何が新しいのか」が見えない論文からは，読者も査読者も離れていきます．
%
% ■ 前提知識は切り分けて，先に簡潔にまとめます
% 自分の貢献ではない記法・用語・既存技術の説明は，「準備」などの節にまとめます．
% これは，「どこからが自分の貢献なのか」の境界を読者に示す役割を持ちます．
%
% ■ トップダウンに書きます
% 読者が「この先どこへ向かうのか」を常に見通せるよう，全体像 → 詳細の順で書いてください．
% 細部を読み飛ばしても手法の概要が掴めるのが理想です．
%
% ■ 設計上の選択には，必ず「なぜそうしたのか」の理由を添えます
% 「Aを採用した」ではなく，「Bでは〜という問題が生じるため，Aを採用した」と書いてください．
% 理由の書かれていない設計は，読者からは「なんとなくそうした」ようにしか見えません．
%
% ■ 語るべきは「完成した手法の物語」であって，「そこに至った試行錯誤の経緯」ではありません
% 自分が苦労した順に書きたくなりますが，読者が知りたいのは完成した手法の筋道です．
% 本筋を中断させる細部（証明の詳細，実装上の細かい工夫など）は，付録への退避を検討してください．
%
% ■ 読者が再現できる詳しさと，本質を見失わない抽象度のバランスを取ります
% パラメータの値やデータ構造の詳細は必要ですが，それらに埋もれて「何をしている手法なのか」が
% 見えなくなってはいけません．まず手法の全体像を1段落で述べてから，細部に入ってください．


% ----------------------
% 提案手法
% ----------------------
\section{提案手法}